\chaptertitle{第八章\quad 抽屉原理、染色与不变量}

前面的章节都是关于"有多少种"的计数问题。本章换一个角度——研究"能不能保证"的存在性问题。抽屉原理是这个领域最基础也最强大的工具，而染色法和不变量法则是构造抽屉的两把利器。

\sectiontitle{8.1\quad 第一抽屉原理}

先看三个简单的问题：
\begin{enumerate}
\item $4$个苹果放进$3$个抽屉，是否一定有一个抽屉里至少有$2$个苹果？
\item $13$个人中，是否一定有$2$个人同一个月出生？
\item $367$个人中，是否一定有$2$个人同一天生日？
\end{enumerate}

直觉告诉我们：是的。这就是\textbf{第一抽屉原理}（鸽巢原理）：
\[
\text{把$n+1$个物体放进$n$个抽屉，必有一个抽屉里至少有$2$个物体。}
\]

\begin{example}
证明：任意$5$个整数中，必有$2$个数的差是$4$的倍数。
\end{example}
\begin{solution}
一个数除以$4$，余数有$0,1,2,3$四种可能（$4$个抽屉）。\\
$5$个数放入$4$个余数类，必有两个数余数相同。\\
这两个数的差是$4$的倍数。
\end{solution}

\begin{example}
$1$到$10$中任取$6$个数，必有两个数的和是$11$。
\end{example}
\begin{solution}
把数配对：$(1,10)$、$(2,9)$、$(3,8)$、$(4,7)$、$(5,6)$，共$5$对（$5$个抽屉）。\\
取$6$个数，根据抽屉原理，必有两个数来自同一对，它们的和为$11$。
\end{solution}

\begin{exercise}
\item 任意$3$个整数中，必有两个数的和是偶数。为什么？
\item $1$到$8$中任取$5$个数，必有两个数的差为$2$。为什么？
\end{exercise}

\sectiontitle{8.2\quad 推广形式}

第一抽屉原理可以推广到更一般的情形：
\[
\text{把$m$个物体放进$n$个抽屉（$m>n$），必有一个抽屉里至少有$\lceil m/n\rceil$个物体。}
\]
其中$\lceil x\rceil$表示不小于$x$的最小整数。

\begin{example}
把$100$个苹果放进$9$个抽屉，必有一个抽屉至少有多少个苹果？
\end{example}
\begin{solution}
$\lceil 100/9\rceil = \lceil 11.11\ldots\rceil = 12$个。
\end{solution}

\begin{example}
一个班有$45$人，至少有多少个人同一个月出生？
\end{example}
\begin{solution}
$45$人放进$12$个月：$\lceil 45/12\rceil = \lceil 3.75\rceil = 4$人。
\end{solution}

\textbf{平均值形式}：如果$n$个正数的平均数是$A$，则至少有一个数$\ge A$，也至少有一个数$\le A$。

\begin{example}
某次考试$5$人的平均分是$85$分，是否一定有人的分数$\ge85$？是否一定有人的分数$>85$？
\end{example}
\begin{solution}
一定有分数$\ge85$（如果都小于$85$，平均数也小于$85$）。\\
不一定有分数$>85$（可能全部恰好$85$）。
\end{solution}

\begin{exercise}
\item $101$个苹果放进$10$个抽屉，必有一个抽屉至少有多少个苹果？
\item 一个袋子里有$4$种颜色的球各$10$个，至少取多少个才能保证有$3$个同色？
\end{exercise}

\sectiontitle{8.3\quad 构造抽屉的技巧}

抽屉原理的关键在于"如何构造抽屉"。同一个问题，不同的抽屉构造会导致不同的结论。

三个常用技巧：
\begin{enumerate}
\item \textbf{分区法}：用配对、分组的方式划分抽屉。
\item \textbf{分类法}：按余数、颜色、属性分类。
\item \textbf{染色法}：用颜色标记物体，同色的归入同一个抽屉。
\end{enumerate}

\textbf{染色问题}是构造抽屉的绝佳载体。

\begin{example}
一个$3\times3$的棋盘，用黑、白两种颜色给每个格子染色。证明：无论怎么染，一定存在一个$2\times2$的子棋盘，其中某种颜色至少出现$3$次。
\end{example}
\begin{solution}
把$3\times3$棋盘分成$4$个重叠的$2\times2$子棋盘（左上、右上、左下、右下）。\\
每个$2\times2$子棋盘有$4$格，可能染色方案有$2^4=16$种。\\
如果我们只考虑某种具体构造，不如换个思路：

考虑$3\times3$棋盘共有$9$格，用两种颜色染。由抽屉原理，至少有一色出现$\lceil9/2\rceil=5$次。
\end{solution}

\begin{example}
在边长为$1$的正方形内任意放$5$个点。证明：至少有两个点之间的距离不超过$\dfrac{\sqrt2}{2}$。
\end{example}
\begin{solution}
把正方形分成$4$个边长为$\dfrac12$的小正方形（$4$个抽屉）。\\
$5$个点放入$4$个小正方形，至少有两个点在同一个（或相邻的）小正方形内。\\
一个小正方形的对角线长为$\dfrac{\sqrt2}{2}$，所以这两个点的距离不超过$\dfrac{\sqrt2}{2}$。
\end{solution}

\begin{example}
证明：任意$6$个人中，要么有$3$个人互相认识，要么有$3$个人互不认识。
\end{example}
\begin{solution}
任选一个人$A$，其余$5$人——要么至少有$3$人认识$A$，要么至少有$3$人不认识$A$（抽屉原理）。
\begin{itemize}
\item 如果$A$认识至少$3$人，看这$3$人之间：\\
——如果其中有两人互相认识，则$A$和这两人构成$3$人互相认识。\\
——如果这$3$人两两互不认识，则他们已经构成$3$人互不认识。
\item 如果$A$不认识至少$3$人，同理可得结论。
\end{itemize}
\end{solution}

\sectiontitle{8.4\quad 不变量法}

不变量法：在操作或变化过程中，找到某个保持不变的量，利用这个不变的性质来证明结论。

\begin{example}
一个$8\times8$棋盘去掉两个对角格（一黑一白），能否用$1\times2$的多米诺骨牌不重叠地铺满？
\end{example}
\begin{solution}
棋盘共有$64$格，去掉$2$格还剩下$62$格。$1\times2$骨牌每块盖$2$格，共需$31$块。

关键观察：棋盘是黑白相间的。去掉的两个对角格如果是同色的，则剩下的黑格和白格数量不同。\\
标准的$8\times8$棋盘有$32$黑$32$白。去掉两个同色的对角格后，比如去掉两个黑角，剩下$30$黑$32$白。\\
每块$1\times2$骨牌恰好盖一黑一白，所以能铺满的条件是黑白格数相等。\\
由此可知，去掉两个同色对角格无法铺满（不变量——黑白格数差保持不变）。\\
如果去掉一黑一白，则可以铺满。
\end{solution}

\begin{example}
一个圆桌周围有$n$个座位。甲、乙两人轮流在空座位上坐下，要求不能坐在已有人旁边。谁无法坐下谁输。证明：如果$n$为偶数，后手必胜。
\end{example}
\begin{solution}
不变量——对称性。后手采用"对称策略"：先手坐下后，后手坐在与先手相对的位置（中心对称点）。由于$n$是偶数，中心对称总存在空座。\\
先手坐下后，其"对称"位置也空着（除非被占了，但被占意味着先手坐到了对称位置——这不可能，因为对称位置不同且$n$为偶数——实际上对称位置总是可以坐的，除非先手把那个位置占了）。\\
通过保持对称性，后手永远不会输——最后先手会先找不到空座。
\end{solution}

\begin{example}
墙上写了一个数$123456789$。每次操作可以选择一个两位数，交换它的两个数字。能否通过若干次操作得到$987654321$？
\end{example}
\begin{solution}
交换两位数字不改变这些数字的奇偶性排列？实际上关注不变量：数字和不变。\\
$1+2+3+4+5+6+7+8+9 = 45$。$9+8+7+6+5+4+3+2+1 = 45$。数字和相同，不能判断。

换个不变量：逆序数——一个排列中，$i<j$且$a_i>a_j$的数对个数。
原排列$123456789$的逆序数为$0$。\\
每交换相邻两个数，逆序数改变$1$（增加或减少$1$）。\\
目标排列$987654321$的逆序数为$C_9^2=36$。\\
由于每次交换改变逆序数$1$，可以从$0$一步步变成$36$，所以是可能的。

但这个操作是"交换一个两位数的两个数字"，不是相邻交换。所以需要看是什么操作。
\end{solution}

\begin{exercise}
\item $5\times5$棋盘用$1\times2$骨牌能否铺满？为什么？
\item 一个黑板上写着$1,2,3,\dots,20$。每次操作擦掉两个数$a,b$，写上$a+b+ab$。最后剩下一个数，是多少？（提示：找不变量$(a+1)(b+1)$。）
\end{exercise}

\sectiontitle{本章小结}

抽屉原理说：东西比抽屉多，就必然有抽屉里有至少两个东西。它的推广形式（$\lceil m/n\rceil$）给出了"至少有多少个"的量化结论。构造抽屉的常用技巧有分区、分类和染色。不变量法则在操作变化中寻找不变的性质来证明结论——棋盘染色（黑白格数差）、对称性、逆序数等都是典型的不变量。

\sectiontitle{课后练习}

\subsection*{A组\quad 基础巩固}

\begin{exercise}
\item $31$个人中至少有几个人同一个月出生？
\item 一个袋子里有$3$种颜色的球各$8$个，至少取多少个才能保证有$4$个同色？
\item $1$到$10$中任取$6$个数，必有两数的和是$11$（用抽屉原理解释）。
\item 证明：在$1,2,\dots,10$中任取$6$个数，必有两个数的差是$5$。
\end{exercise}

\subsection*{B组\quad 挑战提升}

\begin{exercise}
\item 在边长为$1$的正六边形内任意放$7$个点，证明至少有两个点的距离不超过$1$。
\item $8\times8$棋盘去掉左上角和右下角（同色格），能否用$1\times2$骨牌铺满？
\item 证明：任意$5$个整数中，必能选出$3$个，它们的和是$3$的倍数。
\item 黑板上有$1,2,\dots,99$。每次操作擦去两个数$a,b$，写上$|a-b|$。最后剩下一个数，它是奇数还是偶数？
\end{exercise}

\vfill
\noindent\rule{\textwidth}{0.4pt}
本章练习答案请参见书末附录。
